% ============================================================================
% CV de Eduardo Apata Tito
% Compilar: latexmk -pdf CV.tex   (o guardar en VSCode con LaTeX Workshop)
%
% Reglas:
%  - Una página. Sin foto ni datos personales como edad, RUT o estado civil.
%  - Orden cronológico inverso: lo más reciente primero en cada sección.
%  - Cada viñeta empieza con un verbo de acción, sin pronombres.
%  - Cuantifica cuando puedas: cifras, porcentajes, usuarios, tiempos.
% ============================================================================

\documentclass[letterpaper,12pt]{article}

\usepackage[T1]{fontenc}
\usepackage[spanish,es-noshorthands]{babel}
\usepackage{stix2}                       % STIX Two (antes: \usepackage{newtxtext}, Times New Roman)
\usepackage[margin=0.5in]{geometry}
\usepackage{enumitem}
\usepackage[hidelinks]{hyperref}

% PDF legible por sistemas ATS (copiar/pegar y lectura automática)
\input{glyphtounicode}
\pdfgentounicode=1    % caracteres con su equivalente Unicode (tildes, ñ)
\pdfinterwordspaceon  % espacios reales entre palabras, no solo separación visual
\hypersetup{
  pdftitle    = {CV - Eduardo Apata Tito},
  pdfauthor   = {Eduardo Apata Tito},
  pdfsubject  = {Desarrollador Fullstack},
  pdfkeywords = {Desarrollo web, Fullstack, Laravel, PHP, JavaScript, TypeScript, Astro, Redes, Ciberseguridad},
  pdflang     = {es-CL}
}

% Sin palabras cortadas con guion al final de línea
\hyphenpenalty=10000
\exhyphenpenalty=10000

\pagestyle{empty}
\setlength{\parindent}{0pt}
\setlength{\parskip}{0pt}
\raggedbottom
\linespread{1.12}

% ----------------------------------------------------------------------------
% Espaciado: todas las distancias verticales del documento salen de aquí
% ----------------------------------------------------------------------------
\newlength{\espacioBajoNombre}    \setlength{\espacioBajoNombre}{2pt}
\newlength{\espacioBajoTitular}   \setlength{\espacioBajoTitular}{2pt}
\newlength{\espacioBajoContactos} \setlength{\espacioBajoContactos}{6pt}
\newlength{\espacioSobreSeccion}  \setlength{\espacioSobreSeccion}{12pt} % antes de cada título
\newlength{\espacioBajoLinea}     \setlength{\espacioBajoLinea}{6pt}     % entre la línea y el contenido
\newlength{\espacioEntreEntradas} \setlength{\espacioEntreEntradas}{6pt}
\newlength{\espacioSobreLogros}   \setlength{\espacioSobreLogros}{1pt}   % entre el cargo y sus viñetas
\newlength{\espacioEntreLogros}   \setlength{\espacioEntreLogros}{1pt}
\newlength{\grosorLinea}          \setlength{\grosorLinea}{0.3pt}

% ----------------------------------------------------------------------------
% Comandos
% ----------------------------------------------------------------------------

% Separador de los datos de contacto
\newcommand{\separador}{\unskip\hspace{0.6em}·\hspace{0.6em}\ignorespaces}

% Línea horizontal de ancho completo. Tras ella, el siguiente renglón se
% separa igual que cualquier par de renglones.
\newcommand{\linea}{\nobreak\hrule height \grosorLinea\relax\prevdepth=0pt}

% \seccion{Título}: mayúsculas, negrita y línea debajo
\newcommand{\seccion}[1]{%
  \par\addvspace{\espacioSobreSeccion}%
  {\bfseries\MakeUppercase{#1}\par}%
  \nobreak\vskip 2pt\linea
  \nobreak\vskip\espacioBajoLinea} % \vskip: se combina con el espacio de lo que sigue

% Renglón con texto a la izquierda y a la derecha
\newcommand{\dosColumnas}[2]{%
  {\rightskip=0pt \parfillskip=0pt \noindent #1\hfill #2\par}}

% \entrada{Organización}{Ciudad, País}{Cargo o título}{Fechas}
\newcommand{\entrada}[4]{%
  \par\addvspace{\espacioEntreEntradas}%
  \dosColumnas{\textbf{#1}}{\textbf{#2}}%
  \dosColumnas{\small #3}{\small\itshape #4}}

% \entradaSimple{Texto}{Fecha}: entrada de un solo renglón
\newcommand{\entradaSimple}[2]{%
  \par\addvspace{\espacioEntreEntradas}%
  \dosColumnas{\small #1}{\small\itshape #2}}

% \categoria{Nombre}{Elementos}: renglón con etiqueta en negrita
\newcommand{\categoria}[2]{{\small\textbf{#1:} #2\par}}

% Viñetas de logros o habilidades
\newenvironment{logros}
  {\begin{itemize}[leftmargin=0.35in, label=\textbullet, before=\small\raggedright,
                   topsep=\espacioSobreLogros, partopsep=0pt,
                   itemsep=\espacioEntreLogros, parsep=0pt]}
  {\end{itemize}}

% ============================================================================
\begin{document}

% ---------------------------------------------------------------- Encabezado
{\centering
{\fontfamily{stix2}\fontsize{28}{32}\selectfont\bfseries\scshape Eduardo Apata Tito\par}
\vspace{\espacioBajoNombre}
{\itshape Desarrollador Fullstack \hspace{0.5em}|\hspace{0.5em} Ingeniero Civil en Computación e Informática\par}
\vspace{\espacioBajoTitular}
{\small
  Arica, Chile \separador
  \href{mailto:eduardotito.jose01@gmail.com}{eduardotito.jose01@gmail.com} \separador
  \href{https://linkedin.com/in/eapatatito}{linkedin.com/in/eapatatito} \separador
  \href{https://github.com/eduardoapatat}{github.com/eduardoapatat}\par}
}
\vspace{\espacioBajoContactos}
\linea
\vspace{\espacioSobreSeccion}

% ------------------------------------------------------------------- Resumen
{\small\itshape\raggedright
  Desarrollador fullstack e Ingeniero Civil en Computación e Informática. Construyo
  sistemas web con Laravel y JavaScript, con experiencia en el sector público en
  análisis de sistemas y arquitectura de la información. Me sigo formando en redes y
  ciberseguridad (Cisco CCNA) y comunico soluciones técnicas a usuarios no técnicos.\par}

% --------------------------------------------------------------- Experiencia
\seccion{Experiencia profesional}

\entrada{Municipalidad de Arica}{Arica, Chile}
{Arquitecto de la Información (práctica profesional)}{Ene 2026 -- Mar 2026}
\begin{logros}
  \item Reorganicé la jerarquía de la información de varios sitios web municipales, incluido el del instituto municipal.
  \item Eliminé contenido duplicado y desactualizado y simplifiqué los menús para llegar a los servicios en menos pasos.
  \item Elaboré una propuesta de reestructuración del sitio web del Parque Centenario, actualmente en evaluación.
  \item Recomendé al equipo web optimizar imágenes, mejorar la accesibilidad y mantener los contenidos al día.
\end{logros}

\entrada{Universidad de Tarapacá}{Arica, Chile}
{Ayudante de Business Intelligence}{Ago 2025 -- Dic 2025}
\begin{logros}
  \item Guié a estudiantes de Inteligencia de Negocios en el ciclo completo de un proyecto de BI.
  \item Orienté el modelado dimensional de data marts y data warehouses, los procesos ETL y los dashboards.
\end{logros}

\entrada{Corporación Nacional de Desarrollo Indígena (CONADI)}{Arica, Chile}
{Analista de Sistemas (práctica)}{Ene 2025 -- Feb 2025}
\begin{logros}
  \item Levanté requerimientos y diseñé un prototipo de sistema de gestión documental adaptado al flujo de trabajo del usuario.
  \item Presenté la propuesta en lenguaje no técnico; fue bien recibida y el usuario validó que respondía a su necesidad.
\end{logros}

% ----------------------------------------------------------------- Educación
\seccion{Educación}

\entrada{Universidad de Tarapacá}{Arica, Chile}
{Ingeniería Civil en Computación e Informática}{Mar 2020 -- Jul 2026}

% ----------------------------------------------------------- Certificaciones
\seccion{Certificaciones}

\entradaSimple{\textbf{CCNA: Introduction to Networks} | Cisco}{Jul 2023}

% --------------------------------------------------------------- Habilidades
\seccion{Habilidades}

\begin{logros}
  \item Liderazgo de proyectos, desde la necesidad del usuario hasta la entrega.
  \item Levantamiento de requerimientos, análisis de sistemas y comunicación con usuarios no técnicos.
  \item Arquitectura de la información, accesibilidad web y gestión de contenidos.
  \item Español nativo.
\end{logros}

% --------------------------------------------------------------- Tecnologías
\seccion{Tecnologías}

\categoria{Desarrollo web}{HTML, CSS, JavaScript, TypeScript, Astro, Tailwind CSS, PHP, Laravel, API REST}
\categoria{Datos y BI}{Modelado dimensional, data warehouses, ETL, dashboards}
\categoria{Redes}{Cisco CCNA}

\end{document}
